\begin{figure}[!htbp]
	\centering
	\resizebox{\textwidth}{!}{%
		\begin{tikzpicture}[
			scale=0.95,
			>=latex,
			trave/.style={line width=1.2pt},
			carico/.style={black, line width=0.4pt},
			forza/.style={black, line width=1pt},
			risultante/.style={gray!65, line width=1.8pt},
			diagonale/.style={gray!50, line width=0.3pt},
			diagT/.style={black, line width=0.8pt},
			diagTb/.style={gray!75, line width=1.0pt},
			diagM/.style={black, line width=0.8pt},
			diagMb/.style={gray!75, line width=1.0pt},
			hatch/.style={gray!40, line width=0.3pt},
			quota/.style={black, line width=0.4pt},
			fibre/.style={gray!70, line width=0.3pt},
			]
			
			% ── Parametri dei vincoli ──
			\def\tw{0.15}    \def\thR{0.28}   \def\rn{0.06}
			\def\rw{0.05}    \def\bw{0.09}    \def\gw{0.22}
			\def\hh{0.10}
			
			% ── Geometria ──
			\def\Lell{4}
			\def\yqi{-0.80}
			
			% ── Parametri del carico distribuito ──
			\def\pgap{0.10}
			\def\php{0.40}
			\pgfmathsetmacro{\pBase}{\pgap}
			\pgfmathsetmacro{\pTop}{\pgap+\php}
			
			% ── Scale dei diagrammi ──
			\def\sT{0.30}
			\pgfmathsetmacro{\Tmax}{(\Lell/2)*\sT}
			\def\sM{0.45}
			\pgfmathsetmacro{\Mmax}{(\Lell*\Lell/8)*\sM}
			
			% =========================================================
			% MACRO RIGA (a): schema strutturale per il pannello N
			% =========================================================
			\newcommand{\PannelloSchemaPartizione}[2]{
				\def\Npart{#1}
				\pgfmathsetmacro{\Lpart}{\Lell/\Npart}
				\pgfmathsetmacro{\Lfreccia}{1.0/\Npart}
				\pgfmathsetmacro{\Rstart}{\pTop+\Lfreccia}
				
				\draw[trave] (0,0) -- (\Lell,0);
				\draw[fibre, dashed] (0, -0.1) -- (\Lell, -0.1);
				
				\draw[thick] (0,0) -- (-\tw,-\thR) -- (\tw,-\thR) -- cycle;
				\draw[thick] (-\gw,-\thR) -- (\gw,-\thR);
				\foreach \x in {-0.16,-0.08,0,0.08,0.16}
				\draw[thin] (\x,-\thR) -- ++(-0.07,-\hh);
				
				\draw[thick] (\Lell,0) -- (\Lell-\tw,-\thR) -- (\Lell+\tw,-\thR) -- cycle;
				\draw[thick] (\Lell-\bw,-\thR-\rw-0.01) circle (\rw);
				\draw[thick] (\Lell+\bw,-\thR-\rw-0.01) circle (\rw);
				\draw[thick] (\Lell-\gw,-\thR-2*\rw-0.04) -- (\Lell+\gw,-\thR-2*\rw-0.04);
				\foreach \x in {-0.16,-0.08,0,0.08,0.16}
				\draw[thin] (\Lell+\x,-\thR-2*\rw-0.04) -- ++(-0.07,-\hh);
				
				\draw[thick, fill=white] (0,0) circle (\rn);
				\draw[thick, fill=white] (\Lell,0) circle (\rn);
				
				\draw[carico] (0,\pBase) rectangle (\Lell,\pTop);
				
				\foreach \kk in {1,...,\numexpr\Npart-1\relax} {
					\pgfmathsetmacro{\xk}{\kk*\Lpart}
					\draw[carico] (\xk, \pBase) -- (\xk, \pTop);
				}
				
				\foreach \kk in {0,...,\numexpr\Npart-1\relax} {
					\pgfmathsetmacro{\xL}{\kk*\Lpart}
					\pgfmathsetmacro{\xR}{(\kk+1)*\Lpart}
					\draw[diagonale] (\xL, \pBase) -- (\xR, \pTop);
					\draw[diagonale] (\xL, \pTop) -- (\xR, \pBase);
				}
				
				\foreach \kk in {0,...,\numexpr\Npart-1\relax} {
					\pgfmathsetmacro{\xb}{(\kk+0.5)*\Lpart}
					\draw[risultante, -{Latex[length=5pt, width=4pt]}]
					(\xb, \Rstart) -- (\xb, 0);
				}
				
				\pgfmathsetmacro{\xbZero}{0.5*\Lpart}
				\node[above, font=\small, gray!50!black] at (\xbZero, \Rstart+0.02) {#2};
				
				\node[above left=-2pt,  font=\small] at (0,0)     {$A$};
				\node[above right=-2pt, font=\small] at (\Lell,0) {$B$};
				
				\foreach \kk in {0,...,\Npart} {
					\pgfmathsetmacro{\xk}{\kk*\Lpart}
					\draw[quota] (\xk,\yqi+0.07) -- (\xk,\yqi-0.07);
				}
				\foreach \kk in {1,...,\Npart} {
					\pgfmathsetmacro{\xL}{(\kk-1)*\Lpart}
					\pgfmathsetmacro{\xR}{\kk*\Lpart}
					\draw[<->, quota] (\xL,\yqi) -- (\xR,\yqi);
				}
				\foreach \kk in {1,...,\Npart} {
					\pgfmathsetmacro{\xMidPart}{(\kk-0.5)*\Lpart}
					\node[above, font=\footnotesize] at (\xMidPart, \yqi-0.08) {$\ell/\Npart$};
				}
			}
			
			% =========================================================
			% MACRO RIGHE (b) e (c): trave appoggiata semplificata
			% =========================================================
			\newcommand{\TraveAppoggiata}{
				\draw[trave] (0,0) -- (\Lell,0);
				\draw[thick] (0,0) -- (-\tw,-\thR) -- (\tw,-\thR) -- cycle;
				\draw[thick] (\Lell,0) -- (\Lell-\tw,-\thR) -- (\Lell+\tw,-\thR) -- cycle;
				\draw[thick, fill=white] (0,0) circle (\rn);
				\draw[thick, fill=white] (\Lell,0) circle (\rn);
				\node[above left=-2pt,  font=\small] at (0,0)     {$A$};
				\node[above right=-2pt, font=\small] at (\Lell,0) {$B$};
			}
			
			% Hatching parabolico (per riga c)
			\newcommand{\HatchingParabola}{
				\foreach \x in {0.10, 0.30, ..., 3.90} {
					\pgfmathsetmacro{\h}{-(\sM/2)*\x*(\Lell-\x)}
					\draw[hatch] (\x, 0) -- (\x, \h);
				}
			}
			
			% Parabola del caso continuo (per riga c)
			\newcommand{\ParabolaNera}{
				\draw[diagM]
				(0, 0)
				\foreach \x in {0.20, 0.40, ..., 3.80} {
					-- ({\x}, {-(\sM/2)*\x*(\Lell-\x)})
				}
				-- (\Lell, 0);
			}
			
			% Hatching triangolare (per riga b)
			\newcommand{\HatchingTriangolo}{
				\foreach \x in {0.10, 0.30, ..., 3.90} {
					\pgfmathsetmacro{\h}{(\Lell/2 - \x)*\sT}
					\draw[hatch] (\x, 0) -- (\x, \h);
				}
			}
			
			% Triangolo lineare del caso continuo (per riga b)
			\newcommand{\TriangoloNero}{
				\draw[diagT] (0,0) -- (0,\Tmax) -- (\Lell,-\Tmax) -- (\Lell,0);
			}
			
			% =========================================================
			% RIGA (a) — Schemi strutturali con partizione (yshift=0)
			% =========================================================
			\begin{scope}[xshift=0cm, yshift=0cm]
				\PannelloSchemaPartizione{2}{$\dfrac{\bar{p}\,\ell}{2}$}
			\end{scope}
			\begin{scope}[xshift=5cm, yshift=0cm]
				\PannelloSchemaPartizione{3}{$\dfrac{\bar{p}\,\ell}{3}$}
			\end{scope}
			\begin{scope}[xshift=10cm, yshift=0cm]
				\PannelloSchemaPartizione{4}{$\dfrac{\bar{p}\,\ell}{4}$}
			\end{scope}
			
			% =========================================================
			% RIGA (b) — Diagrammi del taglio T(z) (yshift=-3.5cm)
			% =========================================================
			% Pannello (b,2)
			\begin{scope}[xshift=0cm, yshift=-2.0cm]
				\TraveAppoggiata
				\HatchingTriangolo
				\draw[diagTb]
				(0,0) -- (0,\Tmax)
				-- (1.0,\Tmax) -- (1.0,0)
				-- (3.0,0) -- (3.0,-\Tmax)
				-- (\Lell,-\Tmax) -- (\Lell,0);
				\TriangoloNero
			\end{scope}
			% Pannello (b,3)
			\begin{scope}[xshift=5cm, yshift=-2.0cm]
				\TraveAppoggiata
				\HatchingTriangolo
				\pgfmathsetmacro{\TII}{\Tmax/3}
				\pgfmathsetmacro{\zbI}{\Lell/6}
				\pgfmathsetmacro{\zbII}{\Lell/2}
				\pgfmathsetmacro{\zbIII}{5*\Lell/6}
				\draw[diagTb]
				(0,0) -- (0,\Tmax)
				-- (\zbI,\Tmax)   -- (\zbI,\TII)
				-- (\zbII,\TII)   -- (\zbII,-\TII)
				-- (\zbIII,-\TII) -- (\zbIII,-\Tmax)
				-- (\Lell,-\Tmax) -- (\Lell,0);
				\TriangoloNero
			\end{scope}
			% Pannello (b,4)
			\begin{scope}[xshift=10cm, yshift=-2.0cm]
				\TraveAppoggiata
				\HatchingTriangolo
				\pgfmathsetmacro{\TII}{\Tmax/2}
				\pgfmathsetmacro{\zbI}{\Lell/8}
				\pgfmathsetmacro{\zbII}{3*\Lell/8}
				\pgfmathsetmacro{\zbIII}{5*\Lell/8}
				\pgfmathsetmacro{\zbIV}{7*\Lell/8}
				\draw[diagTb]
				(0,0) -- (0,\Tmax)
				-- (\zbI,\Tmax)   -- (\zbI,\TII)
				-- (\zbII,\TII)   -- (\zbII,0)
				-- (\zbIII,0)     -- (\zbIII,-\TII)
				-- (\zbIV,-\TII)  -- (\zbIV,-\Tmax)
				-- (\Lell,-\Tmax) -- (\Lell,0);
				\TriangoloNero
			\end{scope}
			
			% =========================================================
			% RIGA (c) — Diagrammi del momento M(z) (yshift=-6.5cm)
			% =========================================================
			% Pannello (c,2)
			\begin{scope}[xshift=0cm, yshift=-3.5cm]
				\TraveAppoggiata
				\HatchingParabola
				\ParabolaNera
				\draw[diagMb]
				(0,0) -- (1.0,-\Mmax) -- (3.0,-\Mmax) -- (\Lell,0);
				\node[above, font=\small] at (0.5*\Lell, -\Mmax+0.02) {$\tfrac{\bar{p}\ell^2}{8}$};
			\end{scope}
			% Pannello (c,3)
			\begin{scope}[xshift=5cm, yshift=-3.5cm]
				\TraveAppoggiata
				\HatchingParabola
				\ParabolaNera
				\pgfmathsetmacro{\MbI}{\Mmax*8/12}
				\pgfmathsetmacro{\MbII}{\Mmax*10/9}
				\pgfmathsetmacro{\zbI}{\Lell/6}
				\pgfmathsetmacro{\zbII}{\Lell/2}
				\pgfmathsetmacro{\zbIII}{5*\Lell/6}
				\draw[diagMb]
				(0,0)
				-- (\zbI,-\MbI)
				-- (\zbII,-\MbII)
				-- (\zbIII,-\MbI)
				-- (\Lell,0);
				\node[above, font=\small] at (0.5*\Lell, -\Mmax+0.02) {$\tfrac{\bar{p}\ell^2}{8}$};
				\node[below, font=\small, gray!50!black] at (0.5*\Lell, -\MbII-0.0) {$\tfrac{5\bar{p}\ell^2}{36}$};
			\end{scope}
			% Pannello (c,4)
			\begin{scope}[xshift=10cm, yshift=-3.5cm]
				\TraveAppoggiata
				\HatchingParabola
				\ParabolaNera
				\pgfmathsetmacro{\MbI}{\Mmax/2}
				\pgfmathsetmacro{\zbI}{\Lell/8}
				\pgfmathsetmacro{\zbII}{3*\Lell/8}
				\pgfmathsetmacro{\zbIII}{5*\Lell/8}
				\pgfmathsetmacro{\zbIV}{7*\Lell/8}
				\draw[diagMb]
				(0,0)
				-- (\zbI,-\MbI)
				-- (\zbII,-\Mmax)
				-- (\zbIII,-\Mmax)
				-- (\zbIV,-\MbI)
				-- (\Lell,0);
				\node[above, font=\small] at (0.5*\Lell, -\Mmax+0.02) {$\tfrac{\bar{p}\ell^2}{8}$};
			\end{scope}
			
			% =========================================================
			% Etichette di colonna in fondo (sotto la riga c)
			% =========================================================
			\node[font=\small\bfseries] at (0.5*\Lell + 0cm, -5.0) {$N=2$};
			\node[font=\small\bfseries] at (0.5*\Lell + 5cm, -5.0) {$N=3$};
			\node[font=\small\bfseries] at (0.5*\Lell + 10cm, -5.0) {$N=4$};
			
		\end{tikzpicture}%
	}
	\caption{Sostituzione tra carico distribuito e
		carichi concentrati per la trave appoggiata,
		in partizioni di $N = 2, 3, 4$ tratti di
		uguale ampiezza (colonne, da sinistra a
		destra). Righe, dall'alto al basso: schema
		strutturale, diagramma del taglio $T(z)$,
		diagramma del momento flettente $M(z)$. In
		nero il caso continuo, in grigio il caso a
		partizione discreta. Nella riga (a) le
		diagonali in ciascuna sotto-area individuano
		il baricentro in cui è applicata la
		risultante $\bar{p}\ell/N$.}
	\label{fig:stat_partizioni_01}
\end{figure}
